\section{Self-sums with prime support in a fixed coset}\label{sec:cosets}

The comparison theorem also covers a prime coset which contains no
identity class. The additional arithmetic issue is that its character
twists can have negative or genuinely complex singular exponents.
Those branches enter the model used for centering, even though they
are smaller than its positive leading part.

Fix an even integer \(m\), put \(G=(\Z/m\Z)^\times\), and let
\(K\le G\), \(a\in G\). Write
\[
 S=aK,\qquad H=\langle K,a\rangle,\qquad
 r=[G:K],\qquad d=\operatorname{ord}_{G/K}(aK),\qquad \rho=1/r.
\]
Define
\[
 \mathcal P_S=\{p:p\nmid m,\ p\bmod m\in S\},\qquad
 C_S=\langle\mathcal P_S\rangle,\qquad f=\1_{C_S}.
\]
Thus \(1\in C_S\). The generated residue group is \(H\), whereas
the prime density is \(\rho\); these are different data. Put
\begin{align}
 c_S&=\frac1{\Gamma(\rho)}\lim_{s\downarrow1}(s-1)^\rho
                        \prod_{p\in\mathcal P_S}(1-p^{-s})^{-1},
                         \label{eq:coset-marginal}\\
 \Gamma_S(N)&=\frac{m}{|H|^2}
       \#\{(u,v)\in H^2:u+v\equiv N\pmod m\},\notag\\
 \mathcal E_S&=\{p:p\nmid m,\ p\bmod m\notin S\},\notag\\
 C_{\mathcal E_S}&=\prod_{p\in\mathcal E_S}
                 \left(1-\frac1{(p-1)^2}\right),\qquad
 J_S(N)=\prod_{\substack{p\mid N\\p\in\mathcal E_S}}
                                      \frac{p-1}{p-2}.\notag
\end{align}
The constant \(c_S\) is positive; an explicit expression is given below.
For \(d>1\), set
\[
 g_S=\rho\bigl(1-\cos(2\pi/d)\bigr),\qquad
 \delta_S=\min\left\{\frac1{8192\varphi(m)},\frac{3g_S}{10}\right\}.
\]
For \(d=1\), use \(g_S=+\infty\) and
\(\delta_S=1/(8192\varphi(m))\).

\begin{theorem}\label{thm:coset}
Under the analytic inputs of Section~\ref{sec:inputs}, uniformly in
all sufficiently large targets with \(\Gamma_S(N)>0\),
\[
 \#\{(x,y)\in C_S^2:x+y=N\}
 =c_S^2\Gamma_S(N)C_{\mathcal E_S}J_S(N)
       \frac{N}{(\log N)^{2-2\rho}}
       \left(1+O((\log N)^{-\delta_S})\right).
\]
If \(\Gamma_S(N)=0\), the count is identically zero.
The condition \(\Gamma_S(N)>0\) is equivalent to solvability in
every finite quotient. Consequently it characterizes eventual
representability. The modulus and coset remain fixed, and the error
is uniform in all target prime valuations.
\end{theorem}

The exponent incorporates both discarded banks and the gap between
the leading and secondary character branches. In particular a fixed
positive exponent works for every fixed coset. For \(d\ge2\),
\(1-\cos(2\pi/d)=2\sin^2(\pi/d)\ge8/d^2\), so one may always
replace \(\delta_S\) by the simpler, weaker
\(1/(8192\varphi(m)^3)\). There is no uniformity in growing \(m\).
The new assertion concerns unfiltered self-sums. A filter on the
whole numerical value need not be preserved by multiplication by a
prime in a nonidentity coset.

\subsection{Every fixed-character branch}

For a character \(\psi\), define
\(F_\psi(s)=\sum_n\psi(n)f(n)n^{-s}\), with all characters inflated
to a common modulus containing \(m\). Fourier expansion of the prime
coset gives
\begin{equation}\label{eq:coset-character-factorization}
 F_\psi(s)=G_\psi(s)
        \prod_{\chi\in K^\perp}L(s,\psi\chi)^{\rho\overline{\chi(a)}}.
\end{equation}
The correction is uniformly absolutely convergent for \(\Re s>1/2\).
Indeed the coefficient of \(p^{-s}\) in the logarithm is
\[
 \rho\psi(p)\sum_{\chi\in K^\perp}\overline{\chi(a)}\chi(p)
                =\psi(p)\1_{p\bmod m\in S};
\]
the remaining logarithmic terms are bounded by a fixed multiple of
\(\sum_p p^{-2\Re s}\). Literal Euler factors include all conductor
deletions. Analytic logarithms continued from the Euler half-plane
specify all complex powers.

A zeta factor occurs precisely for a fixed character
\(\psi\in K^\perp\); its exponent is
\begin{equation}\label{eq:coset-exponents}
                       \alpha_\psi=\rho\psi(a).
\end{equation}
All other twists have no singular branch at one. The exponents
\(\rho\) occur exactly for \(\psi\in H^\perp\). The remaining
exponents are \(\rho\) times nontrivial \(d\)-th roots of unity,
each with the same multiplicity. In particular,
\(\rho-\Re\alpha_\psi\ge g_S\) outside \(H^\perp\).

For the untwisted series, an explicit positive correction is
\[
 \log G_{1,p}(s)=\sum_{j\ge2}
 \frac{\1_{p\bmod m\in S}-\1_{(p\bmod m)^j\in S}}{j p^{js}}
 \qquad(p\nmid m).
\]
Thus
\begin{equation}\label{eq:coset-constant}
 c_S=\frac{(\varphi(m)/m)^\rho G_1(1)}{\Gamma(\rho)}
        \prod_{\substack{\chi\in K^\perp\\\chi\ne1}}
                     L(1,\chi)^{\rho\overline{\chi(a)}}.
\end{equation}
Here \(G_1(1)\) is an absolutely convergent positive Euler product.
Conjugate characters are paired. Continuation along \(s>1\) makes
the expression positive and unambiguous. The ordinary one-variable
asymptotic, and its complex twisted versions, also follow from
\cite[Theorem 1.2 and (1.9)--(1.10)]{delabretechetenenbaum}.
That fixed-function theorem is not used in place of the moving-cutoff
estimate which follows.

\subsection{The arithmetic model and its one-variable means}

If \(K=G\), then \(f\) is the indicator of integers coprime to
\(m\), and we use the exact model \(D=f\). Counting the finite
residue classes on the line proves Theorem~\ref{thm:coset} with
\(O_m(1)\) error. In the model construction below assume
\(K\ne G\), so \(\rho<1\) and the model upper bound tends to zero.

Keep an outer scale \(X\), with \(L=\log X\), and let
\[
 y=e^{L^{2/5}},\quad \mathcal S=\{p:p\le y\},\quad
 V_y=\prod_{p\le y}(1-1/p),\quad E=e^{\sqrt L}.
\]
For \(n=A v\), where \(A=n_{\mathcal S}\), the integer \(v\) is
\(\mathcal S\)-rough. The outer scale and cutoff are held fixed
under multiplication by bank primes.
For \(\psi\in K^\perp\), write
\[
 F_{\psi,R}(s)=F_\psi(s)
        \prod_{\substack{p\le y\\p\in\mathcal P_S}}
                         (1-\psi(p)p^{-s})
          =(s-1)^{-\alpha_\psi}H_{\psi,R}(s).
\]
Since \(\psi(p)=\psi(a)\) on the allowed prime set, fixed-character
Mertens and a disk of radius \(c/\log y\) give
\begin{equation}\label{eq:coset-rough-regular}
 |H_{\psi,R}(1)|\asymp(\log y)^{-\Re\alpha_\psi},\qquad
 \frac{|H_{\psi,R}^{(j)}(1)|}{j!}
       \ll_j|H_{\psi,R}(1)|(C\log y)^j.
\end{equation}
For the derivative bound, the logarithm of each finite Euler product
relative to its value at one is
\(O(|s-1|\sum_{p\le y}(\log p)/p)=O(1)\) on that disk;
the fixed regular factor is analytic and nonzero there.

Choose a sufficiently large fixed depth \(J\), and put
\begin{equation}\label{eq:coset-hankel}
 M_\psi(t)=\sum_{j=0}^J\frac{H_{\psi,R}^{(j)}(1)}{j!}
              \frac{(\log t)^{\alpha_\psi-j-1}}
                   {\Gamma(\alpha_\psi-j)}.
\end{equation}
Write \(M_0=M_1\) for the positive principal branch. On any fixed
power range, \(M_0(t)>0\), and
\begin{equation}\label{eq:coset-secondary-gap}
 |M_\psi(t)|\ll L^{-3g_S/5}M_0(t)
                     \qquad(\psi\notin H^\perp).
\end{equation}
In fact its individual exponent is
\(3(\rho-\Re\alpha_\psi)/5\): the factors \(L\) and \(\log y\)
in \eqref{eq:coset-rough-regular} and \eqref{eq:coset-hankel}
give exactly this ratio. Conjugate branches are conjugate functions.

Every \(\eta\in H^\perp\) is one on the support of \(f\), so
\(F_{\psi\eta,R}=F_{\psi,R}\) and \(M_{\psi\eta}=M_\psi\)
exactly. Define
\begin{equation}\label{eq:coset-model}
 D_J(n)=\frac{f(A)}{V_y}
              \sum_{\psi\in K^\perp}\overline{\psi(v)}M_\psi(v),
 \qquad \xi_J(n)=f(n)-D_J(n).
\end{equation}
Grouping characters modulo \(H^\perp\) factors the bracket by
\(|G/H|\1_{v\bmod m\in H}\). Its leading term is positive and
all other terms satisfy \eqref{eq:coset-secondary-gap}. Therefore
\(D_J\) is real and nonnegative for large \(X\), and
\begin{equation}\label{eq:coset-model-domination}
 D_{0,J}(n)=|G/H|\1_{v\bmod m\in H}\frac{f(A)M_0(v)}{V_y},
 \qquad |D_J-D_{0,J}|\ll L^{-3g_S/5}D_{0,J}.
\end{equation}
For \(d=1\), the difference is zero.

For a precise extension, use the cutoff from
Section~\ref{sec:model}, with the present outer scale \(X\):
set every \(M_\psi(t)=0\) for \(t\le X^{1/4}\), and multiply
\eqref{eq:coset-hankel} by \(\eta_0(\log t/L)\) above that point.
It equals the displayed formula for \(t\ge X^{1/3}\).
The common real cutoff preserves conjugacy,
\(M_{\psi\eta}=M_\psi\), positivity of the combined model,
and \eqref{eq:coset-secondary-gap} on the transition range.
Condition (C3) only uses \(t\gg X^{1/2}/E>X^{1/3}\).
Consequently, on
\(1\le n\le X^3\),
\[
 0\le D_J(n)\ll L^{-3(1-\rho)/5},\qquad |\xi_J(n)|\ll1.
\]
These exponents are independent of the fixed depth. For an allowed
\(p\le y\), multiplication changes \(A\) to \(pA\) and leaves
\(v\) unchanged. Hence, exactly,
\begin{equation}\label{eq:coset-invariance}
 f(pn)=f(n),\qquad D_J(pn)=D_J(n),\qquad \xi_J(pn)=\xi_J(n).
\end{equation}
The characters in \eqref{eq:coset-model} are evaluated at the rough
part \(v\), rather than at the whole integer \(n\).

\begin{lemma}\label{lem:coset-centering}
For every fixed \(B,A_0>0\), a sufficiently large fixed depth \(J\)
gives the recordwise estimate \eqref{eq:comparison-rough-means} for
\(\xi_J\). It holds uniformly for every \(\mathcal S\)-supported
record \(r_0\le E\), every \(Y\in[X^{1/2},X^2]\), bounded periodic
tests of period at most \(L^B\), and the smooth weights permitted
there. In particular it includes every fixed-character contribution.
\end{lemma}
\begin{proof}
First consider an \(\mathcal S\)-rough variable at scale
\(T\in[X^{2/5},X^2]\). Expand a periodic test in characters after
refining its period to \(Q=\operatorname{lcm}(q,m)\). All primes
dividing \(Q\) lie below \(y\). For a character \(\vartheta\bmod Q\),
\eqref{eq:coset-character-factorization} holds literally, with every
small allowed Euler factor removed. A principal constituent occurs
only if \(\vartheta\) is inflated from \(K^\perp\). In that case
the rough series is precisely \(F_{\vartheta,R}\) defined above;
it does not acquire a new conductor-dependent regular part.

Complex powers require a bound on the analytic logarithm, rather
than merely on the modulus of an \(L\)-function. For a principal
character put
\(L^\#(s)=(s-1)L(s)/(s+1)\), and for a nonprincipal character put
\(L^\#=L\). The zero-free input makes \(L^\#\) analytic and
nonzero on \(\Re s>11/12\). Its logarithm is fixed by the Euler
logarithm on \(\Re s=2\), together with the displayed pole-removing
factor. On a fixed disk centered at \(2+it\), contained in that
half-plane, convexity gives
\(\Re\log L^\#\ll\log(Q(2+|t|))\), while the center logarithm is
bounded by the same quantity. Borel--Carath\'eodory on a smaller
fixed disk therefore gives
\begin{equation}\label{eq:coset-analytic-log}
 |\log L^\#(\sigma+it)|\ll\log(Q(2+|t|))
                  \qquad(23/24\le\sigma\le2).
\end{equation}
The disks can have fixed radii because their left boundary remains
strictly to the right of \(11/12\). The consistent center logarithm
is uniformly controlled by its Euler series. Thus every fixed
complex power in \eqref{eq:coset-character-factorization} has
polynomial growth in conductor and height; its imaginary exponent
causes no uncontrolled argument factor.

Use smooth Mellin inversion to height \(\exp(L^{1/10})\), and move
the line to \(1-c/\log y\). This lies inside the preceding strip
for large \(X\). The finite Euler factors cost only \(L^{O(1)}\),
and \eqref{eq:coset-analytic-log} costs \(\exp(O(L^{1/10}))\).
The factor \(T^{s-1}\) supplies \(\exp(-cL^{3/5})\), which dominates
both. Sufficiently many fixed derivatives of the smooth weight
dispose of the Mellin tails. A twist without a principal constituent
therefore has mean \(O_A(TL^{-A})\) for every fixed \(A\).

For \(\vartheta\in K^\perp\), Hankel integration around one gives
exactly the density \(M_\vartheta\) in \eqref{eq:coset-hankel}.
The derivative bound \eqref{eq:coset-rough-regular} gives relative
Taylor remainder \(O_J(L^{-3(J+1)/5})\) for that branch. All branch
magnitudes are at most a fixed multiple of the principal magnitude.
This argument applies without change to negative and complex
\(\alpha_\vartheta\); the reciprocal Gamma factor is the Hankel
integral of \((s-1)^{-\alpha_\vartheta+j}\).

On the model side the unrestricted rough population has series
\(\zeta(s)\prod_{p\le y}(1-p^{-s})\), principal density \(V_y\),
and no other principal branch. Its twist by
\(\vartheta\overline\psi\) is principal exactly when
\(\vartheta=\psi\) on the units modulo \(Q\). Hence
\(V_y^{-1}\sum_{\psi\in K^\perp}\overline\psi(v)M_\psi(v)\)
matches each of the original character means with its exact
coefficient. This proves arbitrary logarithmic centering after
choosing \(J\) beyond the periodic-character and smooth-derivative
costs. A factor \(v^{iu}\), \(|u|\le L^B\), only enlarges those
fixed costs.

Finally, for a fixed record \(r_0\le E\), multiplicativity gives
\(f(r_0v)=f(r_0)f(v)\), and \eqref{eq:coset-model} has the same
factor \(f(r_0)\). The remaining scale is \(T=Y/r_0\), within
the preceding power range for large \(X\). If the record is
forbidden both terms vanish. Otherwise the just-proved estimate
is \(O((Y/r_0)L^{-A_0})\), uniformly in the record and test.
This is the required recordwise assertion, with the depth chosen
after the requested accuracy.
\end{proof}

\subsection{Common banks and the independent sieve check}

Choose one residue \(a_0\in aK\) and take all small and big banks
from primes \(p\equiv a_0\pmod m\), deleting \(p\mid N\).
Their fixed density is \(1/\varphi(m)\), irrespective of \(\rho\).
Equation \eqref{eq:coset-invariance} verifies condition (C2) of
Theorem~\ref{thm:model-comparison}; the bounds above verify (C1),
and Lemma~\ref{lem:coset-centering} verifies (C3).

The progression character frame is compatible with these banks.
If equivalent characters differ by a character \(\theta\bmod m\),
then on a product of a specified \(s\) slots,
\[
 \psi(D)=\theta(a_0)^s\eta(D),\qquad
 \psi(X)=\theta(a_0)^s\eta(X).
\]
The two products have the same slot set, so the same unit phase
occurs in both. The conjugated frame kernel consequently reproduces
\(\psi\) exactly to the existing error. Inequivalent characters
have zero average on the lifted coset, since that average is a unit
phase times their average on its kernel. Siegel--Walfisz and partial
summation give the same logarithmic-cell estimate. Grouping equivalent
characters in a rational-phase expansion cannot increase the total
absolute coefficient mass. Thus there is no new graph construction
or mark-length loss behind this application.

We verify the remaining envelope assumption independently of the
comparison. Put \(\kappa_S=2-2\rho\) and
\(\Sigma_N=NJ_S(N)L^{-\kappa_S}\). This satisfies
\eqref{eq:comparison-scale}, for example with \(\kappa_0=3\), since
\(0\le\kappa_S<2\), \(J_S(N)\ge1\), and the prime-product bound
\(J_S(N)\ll(\log\log(3N))^2\) is uniform in \(N\).
At a bank prime \(p\nmid N\),
the exact two-coordinate valuation law is \eqref{eq:bank-law}.
For coprime bank records \(r_1,r_2\le E\), their divisibilities
select one interval progression of length
\(|\mathcal I|/(r_1r_2)\) on \(x+y=N\).
Exactness excludes two roots at every unused bank prime, and one
remaining root at a prime dividing \(r_1r_2\). Together with
\(1/(r_1r_2)\), these factors are exactly its probability under
the independent bank law. Shared bank-prime divisibility is impossible.

Use the upper Selberg sieve at \(z=N^{1/10}\), with primes up to
\(w=z^{1/16}\). This contains all banks. Every good forbidden prime
has two roots, merging precisely when it divides \(N\). Fixed-coset
prime Mertens gives the following lower bounds for the nonbank
sieve denominator, up to fixed positive constants:
\[
\begin{array}{c|c}
 \text{weights}&\text{sieve denominator}\\ \hline
 f,f&(\log w)^{2(1-\rho)}/J_S(N),\\
 D_J,f&(\log y)^{1-\rho}(\log w)^{1-\rho}/J_S(N),\\
 D_J,D_J&(\log y)^{2(1-\rho)}/J_S(N).
\end{array}
\]
For a model factor use its nonnegative domination by \(2D_{0,J}\),
its bound \(L^{-3(1-\rho)/5}\), and its exclusion of every forbidden
prime at most \(y\). These bounds turn each row into exactly
\(L^{-\kappa_S}J_S(N)\). Fixed residue conditions can be discarded
for this upper estimate. Its CRT endpoint error per retained record
pair is \(O(z^{4+\epsilon})\); summing at most \(E^2\) pairs gives
\(N^{0.45+o(1)}\). The exponent is independent of the model depth.

Rankin's bound with parameter \(c/\log y\) gives
\[
 \#\{n\le N:n_{\mathcal S}>E\}
                 \ll NL^{O(1)}e^{-cL^{1/10}}.
\]
It applies to every actual/model pair because the weights are
bounded. Multiplying the record-stratum estimate by an arbitrary
nonnegative record function bounded by \(L^B\), and summing, proves
\eqref{eq:comparison-envelope}, including its interval-length factor
and its errors. This is condition (C4). Its fixed logarithmic error
exponents do not depend on the eventual density depth.

All four hypotheses of Theorem~\ref{thm:model-comparison} have now
been checked by invariance, one-variable estimates, or a positive
upper sieve. No additive cancellation or residue-energy conclusion
has been included as an arithmetic assumption.

\subsection{The main term and the exact congruence criterion}

\begin{lemma}\label{lem:coset-finite-images}
A residue \(t\bmod M\) is represented by \(C_S\) if and only if
its reduction modulo \(\gcd(M,m)\) is that of some \(h\in H\),
and \(p\nmid t\) at each forbidden prime \(p\mid M\), \(p\nmid m\).
\end{lemma}
\begin{proof}
Necessity follows from prime support. For sufficiency extract an
allowed-prime prefix \(D\) with valuations
\(\min(v_p(t),v_p(M))\) at allowed primes dividing \(M\).
It belongs to \(C_S\), and its residue belongs to \(H\).
The desired unit tail has state \(hD^{-1}\in H\).

The positive monoid generated by \(S\) in the finite group is \(H\),
so choose a nonempty word \(s_1\cdots s_k=hD^{-1}\), with
\(s_i\in S\). An identity state also has a nonempty word, by taking
a positive power of any element of \(S\). Realize all but the last
letter by primes in their reduced classes modulo \(m\), avoiding
\(M\). For the last letter, impose its class modulo \(m\), and
the remaining unit-tail congruences modulo the prime powers of \(M\).
At an allowed prime with prefix valuation \(e<v_p(M)\), this is
the quotient of \(t/p^e\) by the other unit factors modulo
\(p^{v_p(M)-e}\). At a saturated prefix no further congruence is
needed, beyond avoiding that prime. All other good primes require
a unit tail by the hypothesis. At primes dividing \(m\), compatibility
is precisely the prescribed reduction of \(h\).
The resulting CRT class is reduced, and Dirichlet's theorem supplies
a prime in it. Multiplying all chosen primes and \(D\) realizes
the required residue. A single prime would not suffice in general:
the tail state need not itself lie in \(S\).
\end{proof}

\begin{proof}[Proof of Theorem~\ref{thm:coset}]
The full-group case was handled by its exact model. Suppose \(K\ne G\).
Theorem~\ref{thm:model-comparison}, applied using the checks above,
replaces both copies of \(f\) by \(D_J\), at the sparse scale
\(NJ_S(N)L^{-\kappa_S}\), with relative error
\(O(L^{-1/(8192\varphi(m))})\). Its interval version and the
upper sieve handle the endpoint strips.

On the retained records and power bulk,
\eqref{eq:coset-model-domination} replaces \(D_JD_J\) by
\(D_{0,J}D_{0,J}\) with relative error \(O(L^{-3g_S/5})\).
Put
\[
 \alpha_0=\varphi(m)/m,\qquad
 E_-(y)=\prod_{\substack{p\le y\\p\nmid m,\ p\notin\mathcal P_S}}
                        (1-p^{-1}).
\]
Since \(\log v=L+O(\sqrt L+\log L)\), the leading model is
\[
 |G/H|\1_{n\bmod m\in H}\,f(n_{\mathcal S})\lambda_X,
 \qquad \lambda_X=\frac{c_SL^{\rho-1}}{\alpha_0E_-(y)},
\]
with relative error \(O(L^{-1/2})\). Here
\(n_{\mathcal S}\bmod m\in H\), so the rough and whole-value
state tests are equivalent for this leading model.

Let \(R_N\) be the number of pairs counted in \(\Gamma_S(N)\).
The raw fixed-residue probability on the integer line is \(R_N/m\).
Its two projections and marginal normalization give
\[
 \frac{R_N}{m}\frac{|G/H|^2}{\alpha_0^2}
              =\frac{mR_N}{|H|^2}=\Gamma_S(N).
\]
This calculation uses \(H\) for the residue support and \(\rho\)
for the logarithmic density; it never identifies \(\rho\) with
\(|H|/|G|\).
At a good forbidden prime the normalized pair factor is
\[
 \frac{1-(2-\1_{p\mid N})/p}{(1-1/p)^2}.
\]
The fundamental sieve lemma at a fixed power level, as in
\cite[Theorem 11.7]{richert}, therefore evaluates the leading model
as
\[
 c_S^2\Gamma_S(N)C_{\mathcal E_S}J_S(N)
                       \frac{N}{L^{\kappa_S}}.
\]
Its sieve error is smaller than every fixed logarithmic power;
CRT endpoints are \(N^{0.45+o(1)}\). The omitted ordinary Euler
tail is \(O(1/y)\), and its target-prime tail is
\(O(L/(y\log y))\), uniformly in all target valuations.
The density replacement and secondary-branch errors are bounded
by the asserted \(L^{-\delta_S}\). This proves the asymptotic.

For the local assertion, testing modulo \(m\) proves necessity of
\(\Gamma_S(N)>0\). Conversely fix such a root. At primes dividing
\(m\), lift the equation while preserving its unit states.
At every other prime, both coefficients are units and \(p\ge3\).
The equation modulo \(p\) has \(p\) pairs; requiring a unit on
each forbidden coordinate removes at most two. A surviving pair
lifts, and CRT combines the choices. Lemma~\ref{lem:coset-finite-images}
realizes each coordinate by an actual element of \(C_S\).
This proves solvability in every finite quotient. Finally the
positive asymptotic proves eventual global representability.
\end{proof}

\subsection{Two cosets which test the branches}

For \(m=4\), \(K=\{1\}\), \(a=3\), the generated group is
\(H=G\), while \(\rho=1/2\). The exponents are \(1/2,-1/2\).
Let \(C^-_4\) contain the positive integers supported on primes
\(3\pmod4\). With
\(P_3=\prod_{p\equiv3(4)}(1-p^{-2})^{-1}\), its marginal is
\(c_S^2=2P_3/\pi^2\). Indeed
\[
 F(s)^2=(1-2^{-s})\frac{\zeta(s)}{L(s,\chi_{-4})}F(2s).
\]
Here \(\Gamma_S(N)=2\) on every even target. Since
\(\prod_{p\text{ odd}}(1-p^{-2})^{-1}=\pi^2/8\),
\[
 2c_S^2C_{\mathcal E_S}
       =\frac12\prod_{p\equiv1(4)}
                    \left(1-\frac2{p(p-1)}\right)=:A_-.
\]
Thus Theorem~\ref{thm:coset} gives
\[
 \#\{x+y=N:x,y\in C^-_4\}
 = A_-\prod_{\substack{p\mid N\\p\equiv1(4)}}\frac{p-1}{p-2}
       \frac N{\log N}
       \left(1+O((\log N)^{-1/16384})\right)
\]
for every sufficiently large even \(N\). Odd targets have zero
count. The negative branch is needed for centering, although it
does not change this leading constant.

For a genuinely complex example, take \(m=10\), \(K=\{1\}\),
\(a=3\), and let \(C_{3,10}\) be supported on primes
\(3\pmod{10}\). Choose the literal quartic character \(\chi\bmod10\)
with \(\chi(3)=i\). Then \(H=G\), \(\rho=1/4\), \(d=4\), and
the exponents of \(1,\chi,\chi^2,\overline\chi\) are
\[
                       \frac14,\quad\frac i4,\quad-\frac14,
                                      \quad-\frac i4.
\]
The rough model is explicitly real:
\[
 D_J(n)=\frac{f(A)}{V_y}
   \left\{M_1(v)+2\Re\bigl(\overline{\chi(v)}M_\chi(v)\bigr)
                       +\chi(v)^2 M_{\chi^2}(v)\right\}.
\]
The complex terms have relative size \(O(L^{-3/20})\), and the
negative branch has relative size \(O(L^{-3/10})\). They cannot
be omitted from an arbitrarily accurate centered model.
The leading constant is given by \eqref{eq:coset-marginal}; equivalently,
\[
 c_{3,10}=\frac{(2/5)^{1/4}G_1(1)}{\Gamma(1/4)}
       L(1,\chi)^{-i/4}L(1,\chi^2)^{-1/4}
                              L(1,\overline\chi)^{i/4},
\]
with the positive determination and convergent correction specified
in \eqref{eq:coset-constant}. The forbidden good residues are
\(1,7,9\pmod{10}\). For even \(N\),
\[
 \Gamma_S(N)=
 \begin{cases}5/2,&5\mid N,\\15/8,&5\nmid N.\end{cases}
\]
This follows by counting pairs of nonzero residues modulo five:
there are four at a zero target and three otherwise. Consequently
every even target is congruence-admissible, and
\[
 \#\{x+y=N:x,y\in C_{3,10}\}
 =c_{3,10}^2\Gamma_S(N)C_{\mathcal E_S}J_S(N)
       \frac N{(\log N)^{3/2}}
       \left(1+O((\log N)^{-1/32768})\right).
\]
This includes all sufficiently large even targets, uniformly in
their prime valuations. The example tests complex character
contributions without changing the comparison theorem or its inputs.

The common-bank condition is substantive. Two different prime
prescriptions can be disjoint and supply no common allowed prime
class. Theorem~\ref{thm:coset} makes no assertion about arbitrary
pairs of prime cosets.
